\section{Introduction}
The distributional challenge of artificial intelligence is not identical to its production challenge. An economy may become better at producing goods and services while the institutions through which households obtain claims on that production change more slowly. Labour displacement is one possible source of this mismatch; it is not an inevitable implication of every AI application. Recent workplace evidence documents substantial augmentation, while models of more extensive automation admit very different long-run factor-income distributions \citep{brynjolfsson2025,acemoglu2018,trammell2026}. The policy question is how to support an uncertain transition without confusing an increase in productive capacity with an automatic improvement in household access to it.

This paper investigates \emph{contribution-weighted usownership}. The noun formalises a line of work that predates the present macroeconomic model. In 2022, the standards-track ERC-4950 proposal on entangled tokens, co-authored by Mu\~noz and de la Rosa, explicitly listed ``Usowners'' among its applications: users consuming an NFT could become partial owners of that NFT \citep{erc4950}. The present paper generalises that precursor from a paired digital-object use case to a firm-level institution and names the resulting ownership relation \emph{usownership}: \emph{user} and \emph{owner}, with \emph{-ship} denoting the durable relation. A user can progressively become an owner, with the size of the claim depending on the kind and intensity of contribution. The institution begins with predominantly private, investor-owned firms and adds recurring firm-specific equity grants to consumers, engaged users, evaluators, knowledge contributors, prosumers and co-creators. Pooling into a universal social wealth fund remains a comparator or a possible much-later institution, not an initial condition.

The ownership transition is paired with a separate income and monetary layer. Sovereign digital money and a transitional Minimum Vital Income (MVI) provide an income bridge while usownership accumulates. The conjecture is not that monetary issuance creates the real surplus distributed to users. It is that bounded seigniorage can help finance part of the lag between declining labour income and the maturity of distributed capital claims, subject to money demand, inflation and resource constraints.

Why consider this architecture rather than simply improving regulation, taxes or subsidies? An inadequate answer is that regulation changes rules, taxation changes flows and only usownership changes stocks. Property rights are themselves legal rules; a mandatory equity grant is an in-kind redistribution obligation; and taxation can finance asset grants or public investment. Moreover, a fully credible levy and rebate can reproduce the cash flows of an equity distribution. The relevant comparison is therefore between \emph{institutional implementations of redistribution}, not between legislation and its absence. Their differences concern beneficiary selection, control rights, durability, liquidity, risk, investment incentives and administrative capacity.

There is also an earlier Internet-of-Value (IoV) research precursor. De la Rosa Esteva's chapter on IoV systemic risk argued that interoperability was necessary for the democratisation of finance and property and anticipated a ``massive migration of value onto the ledgers'', described as a potential \emph{Value Deluge} whose integrity and ownership would need preservation \citep{delarosa2022iov}. That chapter did not yet contain the present contribution score, corporate dilution rule or monetary bridge. Its relevance is the architecture-level premise: value-bearing claims themselves migrate onto programmable networks, so ownership institutions can become digitally executable rather than merely digitally recorded. ERC-4950 subsequently supplied a narrower technical precursor in which linked tokens could support partial ``Usowner'' status \citep{erc4950}. The present contribution connects those earlier value and ownership primitives to structural change in production and distribution.

A second difference is becoming technologically more relevant. The prevailing financial architecture records money, securities and corporate ownership across separate account systems whose transfers require messaging, reconciliation, custodians and settlement procedures. Tokenisation instead represents money and assets as programmable claims whose ownership and transfer logic can be executed on interoperable platforms. The Bank for International Settlements describes tokenised assets as executable objects and argues that programmable platforms can bring central-bank money, commercial-bank money and financial assets into a common settlement environment \citep{bis2022,bis2023,bis2025}. The European Commission explicitly describes DLT and tokenisation as potentially paving the way towards an ``Internet of Value'' \citep{eciov2026}. This transition does not prove that blockchain is necessary or that tokenisation creates economic value by itself. It changes the implementation frontier: repeated small equity awards, vesting, dividend routing, contribution entitlements and payment can in principle be coordinated by machine-executable rules rather than reconstructed across disconnected ledgers.

This technological complement is especially important for usownership because the institution is granular and recurrent. A conventional employee stock plan may issue shares to thousands of workers; a mature digital platform could have millions of heterogeneous users whose contribution weights change continuously. If every small award requires a separate sequence of registry, custodian, payment and reconciliation events, the mechanism can be operationally prohibitive even when its economics are attractive. A programmable asset-and-money infrastructure can reduce that friction. It cannot remove the legal definition of equity, the dilution borne by existing owners, the need to verify contribution, privacy constraints or the macroeconomic incidence of seigniorage.

We make the economic comparison explicit. First, we derive a cash-flow replication result: given the same distributable profits, beneficiaries and investment response, an appropriately specified dividend levy exactly reproduces participant income from usownership without transferring legal shares. Second, we quantify the ownership-transition and financing constraints jointly. A rising aggregate user share is insufficient when displaced households do not receive grants, and increasing a holding levy eventually erodes its own monetary base. Third, we separate the contribution of equity diffusion from that of monetary redesign using matched counterfactuals rather than attributing a hybrid's whole effect to seigniorage. Fourth, we formalise how programmable monetary and asset infrastructure can reduce the operational cost of repeated ownership allocation without changing the underlying resource transfer.

The numerical laboratory has heterogeneous households, three productive activities, retained-earnings investment, an automation-linked issue rule, endogenous transaction demand and an explicit resource ceiling. It is a closed, post-reform mechanism study, not an estimated model of the European Union. The original design comprises 450 macroeconomic runs; the comparative extension adds 27 policy paths and 72 controlled robustness runs. Exact accounting identities and matched-payout tests make the exercise reproducible. Official European data provide context and units, not an empirical estimate of the behavioural parameters.

The results support a narrower, more useful proposition than unconditional superiority. In the medium-automation path, adding user equity to the universal-transfer comparator reduces year-50 MVI from 19.03\% to 10.57\% of GDP; the additional holding-levy architecture lowers it to 10.16\%. User-held equity reaches 29.67\%. Yet a targeted transfer without new equity costs 7.07\% of GDP, and a matched dividend levy generates exactly the same household cash outcomes as equity under the replication assumptions. The combined ownership and targeted-protection case reduces MVI to 1.91\% while retaining 18.75\% of GDP in residual conventional taxation. These are conditional model outcomes, not forecasts or estimates of an optimal policy.

The strongest case for usownership is therefore not that it makes fiscal policy obsolete. It is that it may institutionalise durable participant claims before and beyond a particular transfer programme, preserve a direct link between participation and the firm, and become materially easier to administer as economic value itself migrates onto programmable digital rails. Its superiority in governance, incentives and implementation cost must be established empirically. Monetary finance remains a bounded auxiliary instrument, not the source of redistributed real surplus. The analysis positions usownership as a candidate component of structural-transition policy alongside public services, competition policy, adaptive taxation and a universal protection floor.
